% Część II, Text Embeddings: uśrednione wektory Word2Vec. Liczby: makra \teW... z generated/te_w2v_macros.tex
% (python -m src.text_analysis, python -m src.text_report). Tekst pisany ręcznie -- przy zmianie wyników sprawdzić wnioski.
\subsection{Uśrednione wektory Word2Vec}\label{sec:te-w2v}

Dodatkowo zbadano osadzenia zdań jako \textbf{średnią wektorów słów Word2Vec} dla dwóch modeli z części Word
Embeddings (sekcja~\ref{sec:we-models}): ogólnego modelu IPI PAN i modelu wytrenowanego na recenzjach PolEmo
(100 wymiarów). Słowa bez wektora pomijano; wektor ma średnio \teWCovIpi\% słów zdania w modelu IPI PAN
i \teWCovPol\% w modelu PolEmo, a żadne zdanie nie zostało bez wektora. Analizę wykonano tymi samymi funkcjami, parami
kontrolnymi i miarami co dla fastText i USE (\texttt{src/text\_analysis.py}), więc tabele porównują pięć modeli.

\begin{figure}[H]
\centering
\includegraphics[width=0.9\textwidth]{tsne_w2v.png}
\caption{t-SNE osadzeń zdań: uśrednione Word2Vec IPI PAN; kolor wg etykiety uzgodnionej (lewy) i korpusu (prawy).}
\label{fig:tsne-w2v}
\end{figure}

\begin{figure}[H]
\centering
\includegraphics[width=0.9\textwidth]{tsne_w2v_polemo.png}
\caption{t-SNE osadzeń zdań: uśrednione Word2Vec PolEmo.}
\label{fig:tsne-w2v-polemo}
\end{figure}

Na wykresach t-SNE (rys.~\ref{fig:tsne-w2v} i~\ref{fig:tsne-w2v-polemo}) klasy są wymieszane, tak jak dla fastText
i USE; w modelu PolEmo widać jedynie skupisko części zdań NEU.

\begin{table}[H]
\centering\small
\begin{tabular}{lrrrrr}
\toprule
Kategoria par & W2V IPI PAN & W2V PolEmo & fastText & USE & TF-IDF \\
\midrule
parafraza & 0{,}710 & 0{,}700 & 0{,}746 & 0{,}492 & 0{,}376 \\
przeciwny sens, podobne słowa & 0{,}882 & 0{,}877 & 0{,}859 & 0{,}794 & 0{,}744 \\
przeciwny sens, te same słowa & 0{,}950 & 0{,}927 & 0{,}892 & 0{,}924 & 0{,}770 \\
ten sam sens, różna długość & 0{,}578 & 0{,}658 & 0{,}604 & 0{,}191 & 0{,}389 \\
kontrola: niepowiązane & 0{,}428 & 0{,}431 & 0{,}506 & 0{,}111 & 0{,}004 \\
\bottomrule
\end{tabular}

\caption{Pary kontrolne (tab.~\ref{tab:probe-means}): średnie podobieństwo kosinusowe dla uśrednionego Word2Vec
i modeli z poprzednich podrozdziałów.}
\label{tab:te-w2v-probes}
\end{table}

\textbf{1. Zdania uznane za podobne, a różne znaczeniowo.} Tak, tak samo jak w fastText. Pary o przeciwnym sensie
i podobnych słowach są bliższe niż parafrazy (\teWOppIpi{} wobec \teWParaIpi{} w modelu IPI PAN,
tab.~\ref{tab:te-w2v-probes}); „Obsługa była bardzo miła'' i „\ldots niemiła'' mają \teWMilaIpi. W korpusie wśród
1\% najbardziej podobnych par zdań pary POZ--NEG stanowią \teWTopOppIpi\% (IPI PAN) i \teWTopOppPol\% (PolEmo),
czyli więcej niż wśród wszystkich par (\teWAllOpp\%), a pary o tej samej etykiecie -- tyle co przypadkiem
(\teWTopSameIpi\% wobec \teWAllSame\%). Najbliższe pary to długie opinie o lekarzach: o podobieństwie decyduje temat,
a negacja to tylko jedno słowo w średniej.

\textbf{2. Zdania podobne znaczeniowo, a odległe, bo użyto innych słów.} Częściowo. Parafrazy (\teWParaIpi) są
wyraźnie bliżej niż zdania niepowiązane (\teWUnrelIpi), bo Word2Vec zna synonimy (\emph{miła}, \emph{uprzejma}), ale
podobieństwo rośnie ze wspólnymi słowami: korelacja Spearmana z miarą Jaccarda zbiorów słów wynosi \teWJacIpi{}
(IPI PAN) i \teWJacPol{} (PolEmo), a pary bez wspólnych słów mają średnio \teWNoSharedIpi, wobec \teWSharedIpi{} dla
par o Jaccardzie powyżej 0{,}1. Np.\ „Ogólnie Czystość i pokoje na 10, atmosfera na 3'' i „Czysto, schludnie,
bardzo wygodne łóżka'' są w modelu IPI PAN daleko od siebie (\emph{czystość} i \emph{czysto} to dla niego różne
słowa), a w modelu PolEmo blisko.

\textbf{3. Zdania o różnej długości i podobnym sensie.} Nie w pełni: krótkie zdanie i jego długa wersja mają
\teWLenIpi{} (IPI PAN) i \teWLenPol{} (PolEmo), mniej niż parafrazy; np.\ „Polecam .'' i długa pochwała hotelu mają
\teWPolecamIpi{} i \teWPolecamPol. Dodatkowe słowa długiego zdania (hotel, rodzina, tydzień) odciągają jego wektor
od oceny.

\begin{table}[H]
\centering\small
\begin{tabular}{llrrrr}
\toprule
Model & Grupa (tokeny) & $n$ & norma & cos do centroidu & śr. cos do innych \\
\midrule
Word2Vec IPI PAN & (4.999, 11.333] & 67 & 2{,}13 & 0{,}747 & 0{,}612 \\
 & (11.333, 20.0] & 70 & 1{,}93 & 0{,}836 & 0{,}686 \\
 & (20.0, 36.0] & 63 & 1{,}81 & 0{,}885 & 0{,}726 \\
 & \emph{Spearman} (tokeny vs. śr. cos) &  &  &  & 0{,}774 \\
\midrule
Word2Vec PolEmo & (4.999, 11.333] & 67 & 2{,}35 & 0{,}718 & 0{,}563 \\
 & (11.333, 20.0] & 70 & 2{,}05 & 0{,}800 & 0{,}628 \\
 & (20.0, 36.0] & 63 & 1{,}91 & 0{,}847 & 0{,}665 \\
 & \emph{Spearman} (tokeny vs. śr. cos) &  &  &  & 0{,}757 \\
\bottomrule
\end{tabular}

\caption{Długość zdania a wektor dla uśrednionego Word2Vec: trzy grupy zdań według liczby tokenów (te same co
w analizie fastText i USE), średnia norma surowego wektora, podobieństwo do centroidu korpusu i do pozostałych zdań.}
\label{tab:te-w2v-length}
\end{table}

\textbf{4. Długość zdania a wartości podobieństwa i „wyrazistość'' wektora.} Tak, silniej niż w fastText
(tab.~\ref{tab:te-w2v-length}). Od zdań krótkich do długich norma wektora maleje (\teWNormShortIpi{} $\to$
\teWNormLongIpi{} w IPI PAN), podobieństwo do centroidu korpusu rośnie (\teWCentShortIpi{} $\to$ \teWCentLongIpi),
a korelacja długości ze średnim podobieństwem do innych zdań wynosi \teWSpearmanIpi{} i \teWSpearmanPol{}
(fastText \teWSpearmanFt). Uśrednianie znosi wkłady słów, więc wektor długiego zdania jest mniej wyrazisty i bliski
„przeciętnemu zdaniu'' -- dlatego najbliższe pary w korpusie to głównie długie zdania. Pomogłoby ważenie słów
(rzadsze ważniejsze) i usunięcie wspólnej składowej wszystkich zdań.

\begin{table}[H]
\centering\small
\resizebox{\textwidth}{!}{\begin{tabular}{lrrrrrrrr}
\toprule
Model & accuracy & macro-F1 & kNN-10 & silhouette & recall POZ & recall NEG & recall NEU & recall AMB \\
\midrule
Word2Vec IPI PAN & 0{,}505 & 0{,}310 & 0{,}329 & $-$0{,}065 & 0{,}44 & 0{,}07 & 0{,}89 & 0{,}00 \\
Word2Vec PolEmo & 0{,}540 & 0{,}348 & 0{,}356 & $-$0{,}048 & 0{,}54 & 0{,}11 & 0{,}89 & 0{,}00 \\
\midrule
fastText & 0{,}505 & 0{,}322 & 0{,}321 & $-$0{,}123 & 0{,}33 & 0{,}15 & 0{,}92 & 0{,}00 \\
USE & 0{,}520 & 0{,}330 & 0{,}334 & $-$0{,}013 & 0{,}50 & 0{,}09 & 0{,}88 & 0{,}00 \\
TF-IDF & 0{,}490 & 0{,}318 & 0{,}335 & $-$0{,}008 & 0{,}32 & 0{,}17 & 0{,}88 & 0{,}00 \\
\bottomrule
\end{tabular}
}
\caption{Klasyfikacja wydźwięku zdania na samych osadzeniach (etykiety uzgodnione; ten sam protokół co
w tab.~\ref{tab:classif}, baza większościowa \teWBase).}
\label{tab:te-w2v-classif}
\end{table}

\textbf{5. Czy osadzenia nadają się do zadania?} Nie samodzielnie (tab.~\ref{tab:te-w2v-classif}). Accuracy wynosi
\teWAccIpi{} (IPI PAN) i \teWAccPol{} (PolEmo) przy bazie większościowej \teWBase; klasyfikator wybiera głównie NEU
(czułość \teWRecNeuIpi), NEG rozpoznaje rzadko (\teWRecNegIpi{} i \teWRecNegPol), a AMB wcale. Model PolEmo wypada
najlepiej z pięciu porównanych (macro-F1 \teWFonePol), ale uczono go na tym samym korpusie recenzji, z którego
pochodzą zdania. Uśrednione Word2Vec kodują głównie temat; do rozpoznawania wydźwięku potrzebny jest model
uwzględniający kolejność słów i negację.
